%!TEX root = ../memoria.tex
% =============================================================================
%  Abstract
% =============================================================================

The digital transformation of the past decades has completely redefined the
forms of human communication, yet it has left aside a sizeable group of users
for whom these tools were not designed: older adults. The result is a digital
divide that, far from narrowing, has been deepened by the growing complexity
of everyday applications and that further isolates a population segment which
already exhibits high rates of unwanted loneliness.

This Bachelor's Thesis presents the design, development and evaluation of
\emph{ConectaMayor}, an accessible web application whose purpose is to
facilitate communication between older adults and their relatives within a
private, simple and secure digital space. The application is organized around
five modules ---a reminder agenda, a contact list, one-to-one messaging, a
shared photo gallery and family linking via invitation code--- and provides
three differentiated user roles with permissions tailored to the way each
family member is expected to use the system: the older adult, the family
editor and the read-only family member.

The system has been implemented using Django, SQLite, Bootstrap 5 and standard
JavaScript with AJAX polling. This choice reflects the need to keep the
codebase manageable for an individual development effort and the suitability
of these tools to the project's requirements. The interface has been designed
following the WCAG 2.1 accessibility guidelines at level AA and the principles
of user-centered design, with particular attention to visual aspects
---typography, contrast and interaction area sizes--- and to the
simplification of the language used throughout.

Validation has been carried out through a usability evaluation with two real
older adult participants, following the methodology proposed by the Spanish
Human-Computer Interaction Association (AIPO). The results confirm that the
application meets the simplicity and comprehensibility goals set out at the
beginning of the project, and have made it possible to identify specific
improvements that have been incorporated into the design prior to the final
submission.
